%% LyX 2.4.4 created this file.  For more info, see https://www.lyx.org/.
%% Do not edit unless you really know what you are doing.
\documentclass[english,t]{beamer}
\usepackage{mathptmx}
\usepackage[T1]{fontenc}
\usepackage[latin9]{inputenc}
\usepackage{amssymb}
\usepackage{cancel}

\makeatletter
%%%%%%%%%%%%%%%%%%%%%%%%%%%%%% Textclass specific LaTeX commands.
% this default might be overridden by plain title style
\newcommand\makebeamertitle{\frame{\maketitle}}%
% (ERT) argument for the TOC
\AtBeginDocument{%
  \let\origtableofcontents=\tableofcontents
  \def\tableofcontents{\@ifnextchar[{\origtableofcontents}{\gobbletableofcontents}}
  \def\gobbletableofcontents#1{\origtableofcontents}
}

%%%%%%%%%%%%%%%%%%%%%%%%%%%%%% User specified LaTeX commands.
\usetheme{Warsaw}
% or ...

\setbeamercovered{transparent}
% or whatever (possibly just delete it)
\setbeamercovered{invisible} %makes overlays invisible


%gets rid of navigation symbols
\setbeamertemplate{navigation symbols}{}

\usepackage{multicol}

\usepackage{tikz}
\usetikzlibrary{calc}
\usetikzlibrary{plotmarks}
\usepackage{pgfplots}
\pgfplotsset{compat=newest}

\usepackage{calc}
\usepackage[nomessages]{fp}

\usepackage{advdate}    % Advancing/saving dates
\usepackage{datetime}   % Dates formatting
\usepackage{datenumber} % Counters for dates
\usepackage[super]{nth}

%\usepackage{pgfpages}
%\pgfpagesuselayout{4 on 1}[a4paper, landscape, border shrink=7mm]
%\pgfpageslogicalpageoptions{1}{border code=\pgfusepath{stroke}}
%\pgfpageslogicalpageoptions{2}{border code=\pgfusepath{stroke}}
%\pgfpageslogicalpageoptions{3}{border code=\pgfusepath{stroke}}
%\pgfpageslogicalpageoptions{4}{border code=\pgfusepath{stroke}}
%%%Don't forget to add 'handout'

\makeatother

\usepackage{babel}
\begin{document}
\newcommand{\xmax}{2000}
\newcommand{\xmin}{0}
\newcommand{\xstep}{0} %must be xmin+n
\newcommand{\ymax}{500}
\newcommand{\ymin}{0}
\newcommand{\ystep}{0} %must be ymin+n

\newcounter{countEx} \setcounter{countEx}{0}
\newcounter{countProb} \setcounter{countProb}{0}
\resetcounteronoverlays{countEx} \resetcounteronoverlays{countProb}

\newcommand{\ex}[3]{
	\addtocounter{countEx}{1}
	\only<#1-#2>{\begin{exampleblock} {Example \chNum.\secNum.\arabic{countEx}.} #3	\end{exampleblock}}
}

\newcommand{\prob}[4]{
	\addtocounter{countProb}{1}
	\only<#1-#2>{\begin{exampleblock} {Problem  \chNum.\secNum.\arabic{countProb}.\,#3} #4	\end{exampleblock}}
}

\newcommand{\yline}[4]{
	(#2 - #4)/(#1 - #3)*(x - #1) + #2
}

\beamerdefaultoverlayspecification{<+->}

%%%%Chapter and Section numbers%%%%%
\newcommand{\chNum}{2}
\newcommand{\secNum}{4}
\newcommand{\secTitle}{Multiplication of Matrices}

\title[Section \chNum.\secNum: \secTitle]{Section \chNum.\secNum: \secTitle}

\subtitle{Finite Mathematics~~MTH 132~~Spring 2014}

\author{Dr.\,James Ashe}

\titlegraphic{\includegraphics[height=1.5cm]{/home/jim/JamesAshe/JCSU/images/jcsu_bull}}

\institute{Johnson C. Smith University}

\date{{}}

\makebeamertitle 
\begin{frame}{Scalar Product of a Matrix}

\begin{columns}


\column{4cm}
\begin{itemize}
\item <1->[]$A=\begin{bmatrix}\begin{array}{rrr}
1 & 0 & -3\\
4 & -2 & 1
\end{array}\end{bmatrix}$
\item <1->[]$-5A=$
\end{itemize}

\column{4cm}
\begin{itemize}
\item <2->[]$B=\begin{bmatrix}\begin{array}{rrr}
1 & 0 & 0\\
0 & 1 & 0\\
0 & 0 & 1
\end{array}\end{bmatrix}$
\item <2->[]$0B=$ 
\end{itemize}

\column{4cm}
\begin{itemize}
\item <3->[]$C=\begin{bmatrix}1\\
2\\
4\\
0
\end{bmatrix}$
\item <3->[]$-10C=$
\end{itemize}
\end{columns}
\end{frame}

\begin{frame}{Matrix Multiplication}

\begin{columns}


\column{4.5cm}
\begin{itemize}
\item <2->[]$\begin{bmatrix}\begin{array}{rrr}
2 & 3 & -1\\
4 & 2 & 0
\end{array}\end{bmatrix}\begin{bmatrix}1\\
8\\
6
\end{bmatrix}$
\end{itemize}

\column{6.5cm}
\begin{block}<3->{When can we find \$AB\$?}
When $A$ is a $m\times n$ and $B$ is a $n\times k$ matrix. 
\end{block}

\begin{block}<4->{Find \$A\$ times \$B\$}
If the $i^{\mbox{th}}$ row of $A$ is $\{a_{1},a_{2},\dots,a_{n}\}$
and the $j^{\mbox{th}}$ column of $B$ is $\{b_{1},b_{2},\dots,b_{n}\}$
then the $i\times j$ entry of $AB$ is: $a_{1}b_{1}+a_{2}b_{2}+\cdots+a_{n}b_{n}$.
\end{block}
\end{columns}
\end{frame}

\begin{frame}{}

\ex{1}{}{Find the product}
\begin{columns}


\column{5.5cm}
\begin{itemize}
\item []$\begin{bmatrix}\begin{array}{rrr}
2 & 3 & -1\\
4 & 2 & 0
\end{array}\end{bmatrix}\begin{bmatrix}1 & 0\\
8 & -1\\
6 & 2
\end{bmatrix}$
\end{itemize}

\column{5.5 cm}

\end{columns}
\end{frame}

\begin{frame}{}

\ex{1}{}{Find the product}
\begin{columns}


\column{5.5cm}
\begin{itemize}
\item []$\begin{bmatrix}\begin{array}{rr}
2 & 3\\
4 & 2
\end{array}\end{bmatrix}\begin{bmatrix}1 & 0\\
8 & -1\\
6 & 2
\end{bmatrix}$
\end{itemize}

\column{5.5 cm}

\end{columns}
\end{frame}

\begin{frame}{}

\ex{1}{}{Find the product}
\begin{columns}


\column{5.5cm}
\begin{itemize}
\item []$\begin{bmatrix}\begin{array}{rr}
2 & 3\\
4 & 2
\end{array}\end{bmatrix}\begin{bmatrix}1 & 0\\
0 & 1
\end{bmatrix}$
\end{itemize}

\column{5.5 cm}

\end{columns}
\end{frame}

\begin{frame}{}

\ex{1}{}{Find the product}
\begin{columns}


\column{5.5cm}
\begin{itemize}
\item []$\begin{bmatrix}\begin{array}{rr}
2 & 3\\
4 & 2
\end{array}\end{bmatrix}\begin{bmatrix}0 & 0\\
0 & 0
\end{bmatrix}$
\end{itemize}

\column{5.5 cm}
\end{columns}
\end{frame}

\end{document}
